\section{Computational and Formal Methods}

The integrity of this non-associative framework relies crucially on rigorous formal verification. Unlike associative geometries, octonionic algebras are notoriously prone to manual coordinate errors, particularly concerning sign conventions in scalar triple products and the orientation of idempotents. 

To unequivocally establish these results, the foundational theorems of this paper were machine-verified using the **Lean 4 interactive theorem prover**, leveraging its dependent type theory and kernel-level validation.

\subsection{Resolving the Non-Associative Expansions}
A major challenge in formalizing non-associative coordinate algebras is that arbitrary bracketings (e.g., $X \cdot (Y \cdot Z)$ vs. $(X \cdot Y) \cdot Z$) yield complex, non-trivial expansions of the off-diagonal vector components. In Lean 4, we bypassed brute-force reduction by utilizing the \texttt{ring} and \texttt{fin\_cases} tactics over the $\mathbb{Z}^3$ components.

For instance, the Baryon associator theorem:
\begin{equation}
[N_{\text{upper}}(u_1), N_{\text{upper}}(u_2), N_{\text{upper}}(u_3)] = -(\mathbf{u}_1 \times \mathbf{u}_2) \cdot \mathbf{u}_3 \, \ell
\end{equation}
was proved by defining a strict extensionality lemma (\texttt{zorn\_ext}), which decomposed the 8-dimensional equivalence into four sub-goals (the two scalars $a,b$ and the two vectors $\mathbf{u},\mathbf{v}$). The \texttt{fin\_cases i} tactic was employed to exhaustively evaluate the vector components across the finite dimension $i \in \{0, 1, 2\}$, allowing the \texttt{ring} tactic to securely annihilate the non-associative cross-product residuals down to $[0,0,0]$.

\subsection{Formalizing the Automorphisms}
To verify the $SU(3)$ gauge symmetry, we constructed the generalized vector transformation $U: \mathbb{Z}^3 \to \mathbb{Z}^3$ and mapped it over the Zorn algebra as \texttt{applyGauge(X, U)}. We mathematically proved that if $U$ preserves both the dot product and the cross product (the defining invariants of $SU(3)$), then the transformation commutes exactly with Zorn multiplication:
\begin{equation}
U(X \cdot Y) = U(X) \cdot U(Y)
\end{equation}
The Lean 4 theorem \texttt{SU3\_gauge\_is\_automorphism} proved this by enforcing strict linearity ($U(ax+by) = aU(x)+bU(y)$) and cross-product invariance across the off-diagonal components. This conclusively embeds the strong-force color group as the algebraic stabilizer of the diagonal causal planes.

\subsection{Symbolic Validation of the Information Manifold}
In parallel with the kernel-verified proofs in Lean 4, a symbolic verification engine was constructed using Python's SymPy library. This symbolic engine computationally evaluated the log-generating potential $\Phi(X)$ and confirmed its exact correspondence to the Zorn dual inverse:
\begin{equation}
\nabla \Phi(X) = -(X^\dagger)^{-1}
\end{equation}
By bridging the Lean 4 type-theoretic proofs with the SymPy algebraic outputs, we established a mathematically airtight guarantee that the topological anomalies of the hadronic spectrum directly map to the metriplectic and information-geometric boundary of the event horizon.
